% GENERATED FILE. Edit canonical lessons, then run npm run generate:books.
% canonical-source-hash: fnv1a64:c7885db2514dd63a
% canonical-lessons: PA-C298-kanda, PA-C298-buta, PA-C298-kikkar, PA-C298-sanghna, PA-C298-kumlauna

\chapter{Thorns and Bushes}
\label{ch:pa-a2-thorns-and-bushes}
\hlchaptermodality{298}

\begin{chapteropening}
\textbf{By the end of this chapter:} \emph{I can describe the bushes and trees by a path.}

Complete the last lesson of chapter 298: I can describe the bushes and trees by a path.
\end{chapteropening}

\section[\texorpdfstring{\pa{ਕੰਡਾ} (kaṇḍā)}{ਕੰਡਾ (kaṇḍā)}]{\pa{ਕੰਡਾ} (kaṇḍā) — a thorn}

\label{lesson:PA-C298-kanda}



\begin{quote}
Before the new one: say the Punjabi for blunt, then the Punjabi for to prick.
\end{quote}

\subsection*{You'll want to know: \pa{ਕੰਡਾ}}
\textbf{\pa{ਕੰਡਾ}} — \emph{kaṇḍā} — \textquotedblleft{}a thorn\textquotedblright{}.

\begin{quote}
\pa{ਮੇਰੇ} \pa{ਪੈਰ} \pa{ਵਿੱਚ} \pa{ਕੰਡਾ} \pa{ਹੈ।}
\end{quote}

\begin{quote}
\emph{mere pair vich kaṇḍā hai.}
\end{quote}

\begin{quote}
There is a thorn in my foot.
\end{quote}

A masculine noun.

\begin{cousinweb}
\textbf{\pa{ਕੰਡਾ}} (\emph{kaṇḍā}) is inherited from Sanskrit \emph{kaṇṭa}, \textquotedblleft{}thorn\textquotedblright{}; CDIAL lists Punjabi \emph{kaṇḍā} under it (CDIAL 2668).
\end{cousinweb}

\begin{grammarlens}[title={What you've built}]
One more word for this chapter.
\end{grammarlens}

\subsection*{Your turn}
\begin{itemize}
\raggedright
  \item \emph{Say it:} \emph{kaṇḍā}
  \item \emph{Say it:} \emph{kaṇḍā}, once more
  \item \emph{Say it:} \emph{mere pair vich kaṇḍā hai.} — There is a thorn in my foot.
  \item \emph{Recall:} say the Punjabi for blunt, then the Punjabi for to prick, then say \emph{kaṇḍā} again
\end{itemize}

\begin{tcolorbox}[breakable,colback=teal!4,colframe=teal!35!black,title={Before you move on}]
What does \pa{ਕੰਡਾ} mean? (\textquotedblleft{}A thorn\textquotedblright{}.) Say it once more.
\end{tcolorbox}
\section[\texorpdfstring{\pa{ਬੂਟਾ} (būṭā)}{ਬੂਟਾ (būṭā)}]{\pa{ਬੂਟਾ} (būṭā) — a bush; a plant}

\label{lesson:PA-C298-buta}



\begin{quote}
Before the new one: say the Punjabi for to prick, then the Punjabi for a thorn.
\end{quote}

\subsection*{You'll want to know: \pa{ਬੂਟਾ}}
\textbf{\pa{ਬੂਟਾ}} — \emph{būṭā} — \textquotedblleft{}a bush; a plant\textquotedblright{}.

\begin{quote}
\pa{ਘਰ} \pa{ਦੇ} \pa{ਸਾਹਮਣੇ} \pa{ਇੱਕ} \pa{ਹਰਾ} \pa{ਬੂਟਾ} \pa{ਹੈ।}
\end{quote}

\begin{quote}
\emph{ghar de sāhmaṇe ikk harā būṭā hai.}
\end{quote}

\begin{quote}
There is a green bush in front of the house.
\end{quote}

A masculine noun.

\begin{cousinweb}
\textbf{\pa{ਬੂਟਾ}} (\emph{būṭā}) is inherited from a reconstructed \emph{*būṭa}, \textquotedblleft{}bush, plant\textquotedblright{}; CDIAL lists Punjabi \emph{būṭā}, \textquotedblleft{}bush, flower\textquotedblright{}, under it (CDIAL 9297).
\end{cousinweb}

\begin{grammarlens}[title={What you've built}]
One more word for this chapter.
\end{grammarlens}

\subsection*{Your turn}
\begin{itemize}
\raggedright
  \item \emph{Say it:} \emph{būṭā}
  \item \emph{Say it:} \emph{būṭā}, once more
  \item \emph{Say it:} \emph{ghar de sāhmaṇe ikk harā būṭā hai.} — There is a green bush in front of the house.
  \item \emph{Recall:} say the Punjabi for to prick, then the Punjabi for a thorn, then say \emph{būṭā} again
  \item \emph{Write it:} copy \textbf{\pa{ਬੂਟਾ}} once, checking the whole word
\end{itemize}

\begin{tcolorbox}[breakable,colback=teal!4,colframe=teal!35!black,title={Before you move on}]
What does \pa{ਬੂਟਾ} mean? (\textquotedblleft{}A bush; a plant\textquotedblright{}.) Say it once more.
\end{tcolorbox}
\section[\texorpdfstring{\pa{ਕਿੱਕਰ} (kikkar)}{ਕਿੱਕਰ (kikkar)}]{\pa{ਕਿੱਕਰ} (kikkar) — an acacia tree}

\label{lesson:PA-C298-kikkar}



\begin{quote}
Before the new one: say the Punjabi for a thorn, then the Punjabi for a bush.
\end{quote}

\subsection*{You'll want to know: \pa{ਕਿੱਕਰ}}
\textbf{\pa{ਕਿੱਕਰ}} — \emph{kikkar} — \textquotedblleft{}an acacia tree\textquotedblright{}.

\begin{quote}
\pa{ਕਿੱਕਰ} \pa{ਦੇ} \pa{ਕੰਡੇ} \pa{ਤਿੱਖੇ} \pa{ਹੁੰਦੇ} \pa{ਹਨ।}
\end{quote}

\begin{quote}
\emph{kikkar de kaṇḍe tikkhe hunde han.}
\end{quote}

\begin{quote}
The thorns of the acacia are sharp.
\end{quote}

A feminine noun.

\begin{cousinweb}
\textbf{\pa{ਕਿੱਕਰ}} (\emph{kikkar}) is inherited from a reconstructed \emph{*kikkara}, \textquotedblleft{}Acacia arabica\textquotedblright{}; CDIAL lists Punjabi \emph{kikkar} under it (CDIAL 3151). Its gender varies: CDIAL marks \emph{kikkar} masculine, beside a feminine \emph{kikkarī}; this lesson treats \pa{ਕਿੱਕਰ} as feminine.
\end{cousinweb}

\begin{grammarlens}[title={What you've built}]
One more word for this chapter.
\end{grammarlens}

\subsection*{Your turn}
\begin{itemize}
\raggedright
  \item \emph{Say it:} \emph{kikkar}
  \item \emph{Say it:} \emph{kikkar}, once more
  \item \emph{Say it:} \emph{kikkar de kaṇḍe tikkhe hunde han.} — The thorns of the acacia are sharp.
  \item \emph{Recall:} say the Punjabi for a thorn, then the Punjabi for a bush, then say \emph{kikkar} again
  \item \emph{Write it:} copy \textbf{\pa{ਕਿੱਕਰ}} once, checking the whole word
\end{itemize}

\begin{tcolorbox}[breakable,colback=teal!4,colframe=teal!35!black,title={Before you move on}]
What does \pa{ਕਿੱਕਰ} mean? (\textquotedblleft{}An acacia tree\textquotedblright{}.) Say it once more.
\end{tcolorbox}
\section[\texorpdfstring{\pa{ਸੰਘਣਾ} (saṅghṇā)}{ਸੰਘਣਾ (saṅghṇā)}]{\pa{ਸੰਘਣਾ} (saṅghṇā) — dense; thick with growth}

\label{lesson:PA-C298-sanghna}



\begin{quote}
Before the new one: say the Punjabi for a bush, then the Punjabi for an acacia tree.
\end{quote}

\subsection*{You'll want to know: \pa{ਸੰਘਣਾ}}
\textbf{\pa{ਸੰਘਣਾ}} — \emph{saṅghṇā} — \textquotedblleft{}dense; thick with growth\textquotedblright{}.

\begin{quote}
\pa{ਇਹ} \pa{ਜੰਗਲ} \pa{ਬਹੁਤ} \pa{ਸੰਘਣਾ} \pa{ਹੈ।}
\end{quote}

\begin{quote}
\emph{eh jaṅgal bahut saṅghṇā hai.}
\end{quote}

\begin{quote}
This forest is very dense.
\end{quote}

\begin{cousinweb}
\textbf{\pa{ਸੰਘਣਾ}} (\emph{saṅghṇā}) is inherited from a reconstructed \emph{*saṁghana}, \textquotedblleft{}very dense\textquotedblright{}; CDIAL lists Punjabi \emph{saṅghṇā}, \textquotedblleft{}thick, close together\textquotedblright{}, under it (CDIAL 12856).
\end{cousinweb}

\begin{grammarlens}[title={What you've built}]
One more word for this chapter.
\end{grammarlens}

\subsection*{Your turn}
\begin{itemize}
\raggedright
  \item \emph{Say it:} \emph{saṅghṇā}
  \item \emph{Say it:} \emph{saṅghṇā}, once more
  \item \emph{Say it:} \emph{eh jaṅgal bahut saṅghṇā hai.} — This forest is very dense.
  \item \emph{Recall:} say the Punjabi for a bush, then the Punjabi for an acacia tree, then say \emph{saṅghṇā} again
  \item \emph{Write it:} copy \textbf{\pa{ਸੰਘਣਾ}} once, checking the whole word
\end{itemize}

\begin{tcolorbox}[breakable,colback=teal!4,colframe=teal!35!black,title={Before you move on}]
What does \pa{ਸੰਘਣਾ} mean? (\textquotedblleft{}Dense; thick with growth\textquotedblright{}.) Say it once more.
\end{tcolorbox}
\section[\texorpdfstring{\pa{ਕੁਮਲਾਉਣਾ} (kumlāuṇā)}{ਕੁਮਲਾਉਣਾ (kumlāuṇā)}]{\pa{ਕੁਮਲਾਉਣਾ} (kumlāuṇā) — to wilt; to wither}

\label{lesson:PA-C298-kumlauna}



\begin{quote}
Before the new one: say the Punjabi for an acacia tree, then the Punjabi for dense.
\end{quote}

\subsection*{You'll want to know: \pa{ਕੁਮਲਾਉਣਾ}}
\textbf{\pa{ਕੁਮਲਾਉਣਾ}} — \emph{kumlāuṇā} — \textquotedblleft{}to wilt; to wither\textquotedblright{}.

\begin{quote}
\pa{ਤੇਜ਼} \pa{ਧੁੱਪ} \pa{ਵਿੱਚ} \pa{ਫੁੱਲ} \pa{ਕੁਮਲਾਉਂਦੇ} \pa{ਹਨ।}
\end{quote}

\begin{quote}
\emph{tez dhupp vich phull kumlāunde han.}
\end{quote}

\begin{quote}
Flowers wilt in strong sunshine.
\end{quote}

Past (it …): \pa{ਉਹ} \pa{ਫੁੱਲ} \pa{ਕੁਮਲਾਇਆ}, \emph{oh phull kumlāiā}.

\begin{cousinweb}
\textbf{\pa{ਕੁਮਲਾਉਣਾ}} (\emph{kumlāuṇā}) is inherited from a reconstructed \emph{*kōmh}, \textquotedblleft{}wither\textquotedblright{}; CDIAL lists Punjabi \emph{kumlāuṇā}, \textquotedblleft{}to wither\textquotedblright{}, under it (CDIAL 3524).
\end{cousinweb}

\begin{grammarlens}[title={What you've built}]
One more word for this chapter.
\end{grammarlens}

\subsection*{Your turn}
\begin{itemize}
\raggedright
  \item \emph{Say it:} \emph{kumlāuṇā}
  \item \emph{Say it:} \emph{kumlāuṇā}, once more
  \item \emph{Say it:} \emph{tez dhupp vich phull kumlāunde han.} — Flowers wilt in strong sunshine.
  \item \emph{Recall:} say the Punjabi for an acacia tree, then the Punjabi for dense, then say \emph{kumlāuṇā} again
  \item \emph{Write it:} copy \textbf{\pa{ਕੁਮਲਾਉਣਾ}} once, checking the whole word
\end{itemize}

\begin{tcolorbox}[breakable,colback=teal!4,colframe=teal!35!black,title={Before you move on}]
What does \pa{ਕੁਮਲਾਉਣਾ} mean? (\textquotedblleft{}To wilt; to wither\textquotedblright{}.) Say it once more.
\end{tcolorbox}
